\Subhead{3.3.1}{Гауссовский процесс. Узкополосный стационарный шум}

Процесс гауссовский, если любой конечный набор его отсчётов имеет совместное нормальное распределение. Для вектора $\bm X$ с невырожденной ковариацией $C$
\begin{equation}
 p(\bm x)=\frac{\exp[-(\bm x-\bm m)^{\mathsf T}C^{-1}(\bm x-\bm m)/2]}
 {(2\pi)^{n/2}\sqrt{\det C}}.
\end{equation}
Линейное преобразование сохраняет гауссовость; некоррелированные компоненты совместно гауссовского вектора независимы. Для негауссовских величин это неверно. Теорема о спариваниях сводит чётные моменты центрированного гауссовского вектора к сумме произведений попарных ковариаций.

Узкополосный шум около $\omega_0$ представим
\begin{equation}
 x(t)=\operatorname{Re}\{z(t)\e^{-\ii\omega_0t}\},\qquad z=X+\ii Y=R\e^{\ii\phi},
\end{equation}
где квадратуры меняются медленно относительно периода. Для кругового комплексного гауссовского шума $X,Y$ независимы, имеют нулевые средние и одинаковую дисперсию $\sigma^2$. Переход к полярным координатам с якобианом $R$ даёт
\begin{equation}
 p_R(r)=\frac r{\sigma^2}\e^{-r^2/(2\sigma^2)},\ r\geq0,
 \qquad p_\phi=\frac1{2\pi}.
\end{equation}
Огибающая имеет распределение Рэлея, интенсивность $R^2$ — экспоненциальное с $\langle R^2\rangle=2\sigma^2$. При добавлении когерентного сигнала огибающая становится райсовской; равномерность фазы исчезает. Физический источник гауссовости — сумма многих слабых независимых вкладов при условиях центральной предельной теоремы, а не сама узкополосность.
